\scanpage{62}
\mspage{5}
% archive foliation: 64
les propriétés de régularité de $C$ en tant que
cercle \inserted{``intérieur''} dans la surface \deleted{sans bord}\inserted{conforme} $X'$ déduite
par bouchage du trou. D'autre part,
le choix d'un point de $D\setminus C$ équivaut
à celui d'une structure torique
sur $C$ compatible avec la structure de
droite projective réelle (i.e. une réduction
du groupe structural \inserted{$(\unsure{\operatorname{GP}(1,\mathbb R)}\simeq)$}
$\operatorname{GL}(2,\mathbb R)/\mathbb R^*$
au sous-groupe cpt maximal $\operatorname O(2)/\pm1$)
--- et cela équivaut aussi à la donnée
d'une structure \unsure{torique|conique} de $D$ sur \inserted{le bord}
$\delta D=C$. Notons que l'on aura alors une
inclusion canonique
\[
D\subset T_{X',a}
\]
où $a$ est le centre choisi de $D$. On
va expliciter \deleted{ainsi} ces diverses
relations ainsi

\uline{Prop.} Équivalence entre \deleted{la cat. des
surfaces conf.}
